\section{Design}

\subsection{System Overview}
\label{sec:design:overview}

Figure~\ref{fig:arch-overview} shows the architecture of \SystemName. The key
idea is to decouple \emph{kernel-enforced confinement} from \emph{semantic
release authorization}. \SystemName treats an agent's permission declaration as
untrusted input: a deterministic SK-IR compiler first translates it into a
bounded, kernel-enforceable upper bound, and the runtime installs the resulting
policy before transferring control to untrusted tool code. The tool then runs
speculatively under this policy while a trusted semantic validator concurrently
decides whether the resulting outputs and effects may be released. This reduces
authorization latency without creating an unconstrained execution window.

\paragraph{Architecture and workflow.}
A tool invocation proceeds in four steps. First, the agent submits a tool
request and an untrusted permission declaration. Second, \SystemName compiles
the declaration into SK-IR, rejecting requests that cannot be safely represented
without expanding authority. The compiled SK-IR is materialized as an
invocation-scoped eBPF/LSM policy, represented by a short-lived \emph{ghost}
object and bound to the tool and its descendants. Third, the tool executes under
kernel confinement while \textsc{ShadowFS} records tentative filesystem effects
and buffers agent-visible outputs, including stdout, stderr, return values, and
structured tool responses. In parallel, the semantic validator evaluates the
task context, accepted permissions, and pending release request. Finally, if
authorization succeeds, \SystemName commits the approved filesystem state and
releases the approved outputs or external effects. If authorization is denied,
times out, or fails before approval, \SystemName discards speculative state,
suppresses buffered outputs, and revokes the corresponding ghost policy.

Operations whose effects cannot be undone after release, such as network
transmissions, API calls, and interactions with remote services, are treated as
release boundaries. They are withheld or routed through mediated release paths
until semantic approval is available; \SystemName does not rely on rolling back
externally observed effects after the fact.

\paragraph{Security invariants.}
This workflow maintains three invariants.

\begin{enumerate}[leftmargin=1.2em]
\item \textbf{Confinement before untrusted execution.}
Kernel enforcement is installed before untrusted tool code runs, so the tool
never executes with ambient host authority.

\item \textbf{Invocation-scoped least privilege.}
Runtime authority is derived only from SK-IR accepted by the deterministic
compiler and is bound to a single invocation. Permissions do not accumulate
across tool calls and are revoked when the invocation completes, fails, times
out, is discarded, or its lease expires.

\item \textbf{Fail-closed release.}
Agent-visible outputs, committed filesystem changes, and externally visible
effects are released only after semantic approval. Denials, timeouts, and
pre-approval failures discard speculative state and revoke the associated ghost
policy.
\end{enumerate}

\subsection{SK-IR: Semantic-to-Kernel Translation}
\label{sec:design:ir}

The challenge in enforcing agent permissions is not only kernel enforcement,
but translating high-level, untrusted permission declarations into
kernel-enforceable constraints in a deterministic and auditable way. Directly
allowing an LLM to generate low-level security policies would place semantic
reasoning on the trusted enforcement path and make the resulting behavior
difficult to verify. \SystemName therefore introduces the Semantic-to-Kernel
Intermediate Representation (SK-IR), a rule-based translation layer between
agent-declared permissions and kernel enforcement.

SK-IR does not decide whether a requested permission is semantically justified.
It compiles the agent's declaration into a precise, auditable upper bound on
runtime authority. The semantic validator separately decides whether the
accepted authority and the resulting pending effects are appropriate for the
task context and may be released.

SK-IR is designed around three goals:

\begin{enumerate}[leftmargin=1.2em]
\item \textbf{Determinism.}
Translation is deterministic with respect to its explicit inputs, including the
permission declaration, compiler version, host policy configuration, resolver
snapshot, and invocation namespace state. It does not depend on LLM behavior.

\item \textbf{Least-privilege preservation.}
Translation may validate, normalize, and narrow requested permissions, but it
must not expand authority. Declarations that cannot be compiled while
preserving this narrowing relationship fail closed.

\item \textbf{Auditability.}
Each deployed enforcement rule can be traced back to the declaration and
compiler decisions that produced it.
\end{enumerate}

\subsubsection{Hard-Boundary Validation}
\label{sec:design:ir:hard}

Before translation, each invocation passes through a hard-boundary validator
that performs syntactic and schema-level checks, including tool registration,
parameter validation, and permission-manifest integrity checks. Requests that
reference unknown tools, contain unsupported parameters, or violate schema
constraints are rejected before policy generation.

Unlike semantic authorization, these checks require no inference. They provide a
deterministic first line of defense against malformed or structurally invalid
requests.

\subsubsection{Multi-Level Translation}
\label{sec:design:ir:translation}

SK-IR defines translation as a partial deterministic function:

\[
\mathcal{T} : \textit{IntentIR} \rightharpoonup \textit{EnforcementIR}.
\]

\noindent The function is partial because not every permission declaration can
be safely represented as kernel-enforceable authority. For accepted
declarations, the compiler maintains an authority-narrowing invariant: under
the compiler-defined authority ordering, the runtime authority represented by
$\mathcal{T}(d)$ is no broader than the authority declared in $d$. If this
relationship cannot be established, translation fails closed and no policy is
emitted.

\paragraph{IntentIR.}
The entry point is a typed representation of the requested operation. Let
$\mathcal{I}$ denote the closed set of supported intent types, such as file
read, file write, network connect, and process execution. Each intent
$d = (i, \pi)$ pairs an intent type $i \in \mathcal{I}$ with a per-type
parameter record $\pi$ drawn from a fixed schema. Free-form policy fragments are
excluded by construction: only parameter keys in a per-type whitelist $K_i$ are
accepted.

\paragraph{ConstraintIR.}
The extraction pass transforms an intent into explicit resource-to-predicate
bindings. It dispatches on the intent discriminator and applies a resource
extraction function $\mathcal{E}_i$ specific to each supported intent type:

\[
\textsc{Extract}(i, \pi) =
\bigl\langle
  \mathcal{E}_i(\pi),\;
  \textit{ops}(i),\;
  \textsc{Allow}
\bigr\rangle .
\]

\noindent The function $\mathcal{E}_i(\pi)$ yields a set of bound resources
$\{(r_j, P_j)\}$, where each typed resource identifier $r_j$ is paired with a
conjunction of constraint predicates $P_j$. For filesystem intents, the trusted
runtime canonicalizes resources with respect to the invocation namespace and
emits predicates over approved namespace roots, mount identifiers, path
components, or stable object identifiers where available. Ambiguous cases,
including unsafe symbolic-link traversal, unauthorized mount crossing, and
sensitive pseudo-paths such as uncontrolled \texttt{/proc} aliases, are rejected
rather than approximated with broader access. For network intents, endpoint
descriptions are converted into predicates over address family, address or CIDR
range, port, and protocol. Capability requests are recorded as symbolic
capability sets.

\paragraph{EnforcementIR.}
The lowering pass rewrites each constraint predicate into a kernel-checkable
form using a fixed set of rewrite rules:

\begin{align}
&\ell(\textit{Endpoint}(\textit{host}, \textit{port}, \textit{proto}))
  \notag\\
&\quad= \textit{AddrSet}\bigl(
    \textsc{Resolve}(\textit{host}),
    \textit{port},
    \textit{proto}
  \bigr)
  \label{eq:endpoint} \\[3pt]
&\ell(\textit{CapSet}(\{c_1, \ldots, c_k\}))
  \notag\\
&\quad= \textit{CapBitmask}\!\left(
    \textstyle\bigvee_{j=1}^{k} 2^{\,\textit{bit}(c_j)}
  \right)
  \label{eq:cap} \\[3pt]
&\ell(\textit{SockFamily}(\{f_1, \ldots, f_m\}))
  \notag\\
&\quad= \textit{FamilyBitmask}\!\left(
    \textstyle\bigvee_{j=1}^{m} 2^{\,\textit{bit}(f_j)}
  \right)
  \label{eq:sock} \\[3pt]
&\ell(\textit{Group}(\textit{op}, \overline{p}))
  \notag\\
&\quad= \textit{Group}(\textit{op},\, \overline{\ell(p)})
  \label{eq:group}
\end{align}

\noindent Rule~\eqref{eq:endpoint} resolves host predicates using a trusted
resolver at compilation time and seals the result as concrete address
predicates with resolver metadata and TTL bounds. Resolution failure, unstable
resolution, or policy-incompatible results cause fail-closed rejection. Single
host addresses are represented as /32 IPv4 or /128 IPv6 masks. Rules
\eqref{eq:cap} and~\eqref{eq:sock} compile symbolic name sets into bitmasks
using static lookup tables that mirror the kernel's \texttt{CAP\_*} and
\texttt{AF\_*} definitions. Rule~\eqref{eq:group} applies lowering recursively
to composite predicates. Predicates that are already kernel-checkable are
preserved; unsupported predicates are rejected rather than passed through.

After predicate lowering, a static hook-selection function chooses a sufficient
LSM hook set for the intent type and lowered predicate set:

\[
H = \textsc{HookSet}(i,\; \{\ell(p) \mid p \in \textstyle\bigcup_j P_j\}) .
\]

\noindent The compiler emits hook-indexed rule entries for the selected hook
classes. The resulting \textit{EnforcementIR} is sealed with default-deny
semantics for managed invocations: operations within the modeled surface that
are not explicitly covered are rejected at the LSM hook level.

Because each pass is a pure function of its typed inputs and explicit host
state, requiring no shared mutable state, randomness, or model inference, the
pipeline is reproducible given the same compiler version, host policy
configuration, resolver snapshot, and invocation namespace state.

\subsubsection{Policy Canonicalization}
\label{sec:design:ir:normalize}

Real-world permissions often contain composite conditions. SK-IR represents
these conditions as logical predicates over resources and access modes, then
canonicalizes them into a normal form for kernel evaluation. The canonical form
maps directly to rule structures stored in eBPF maps, enabling efficient
runtime checks while preserving the intended logical constraints.

Canonicalization performs only semantics-preserving rewrites, such as sorting
commutative predicate groups, eliminating duplicates, merging adjacent address
ranges when doing so does not expand authority, and rejecting contradictory or
unsupported combinations. If canonicalization would exceed configured rule-count
or map-size bounds, \SystemName rejects the request instead of emitting a
partial or overbroad policy.

\subsection{Kernel-Level Enforcement via eBPF/LSM}
\label{sec:design:ebpf}

The output of SK-IR is an \texttt{EnforcementIR} that specifies the
kernel-enforceable upper bound for a tool invocation. The enforcement layer is
responsible for installing this bound before untrusted tool code runs, binding
it to the corresponding invocation, and revoking it when the invocation
terminates or its authorization lease expires.

\SystemName realizes \texttt{EnforcementIR} as eBPF programs attached to Linux
Security Module (LSM) hooks. The runtime places each tool invocation in a fresh
cgroup v2 context, and LSM hooks use this context to select a short-lived
\emph{ghost} policy. This design keeps low-level access-control decisions in
the kernel, while policy construction, provenance tracking, and semantic
release authorization remain in the trusted runtime.

\subsubsection{Invocation-Scoped Policy Enforcement}
\label{sec:design:ebpf:cgroup}

Authorization state must be scoped to a single tool invocation. Process
identifiers are unsuitable because they can be recycled and do not naturally
represent descendant processes created by a tool. Instead, \SystemName binds
each invocation to a fresh cgroup v2 subtree.

Before transferring control to untrusted tool code, the trusted runtime creates
the cgroup, installs the ghost policy state, sanitizes inherited file
descriptors, and places the tool process in the cgroup. Forked and
\texttt{execve}'d descendants inherit the same cgroup membership and therefore
remain subject to the same policy. At each covered LSM hook, the eBPF program
obtains the current cgroup identifier using
\texttt{bpf\_get\_current\_cgroup\_id()} and performs a bounded lookup for the
active ghost policy.

For cgroups managed by \SystemName, missing, expired, or malformed ghost state
is treated as a denial. Confined processes are also prevented from migrating out
of their assigned cgroup, for example by denying access to cgroup control files
and excluding cgroup management interfaces from the confined namespace. This
binds authorization state to the invocation context and prevents permissions
from being reused across concurrent tool invocations.

\subsubsection{LSM-Based Policy Enforcement}
\label{sec:design:ebpf:hooks}

\SystemName enforces the resource classes modeled by SK-IR through a fixed set
of LSM hook classes. Requests requiring operations outside this modeled surface
are rejected during compilation. Within the modeled surface, unmatched runtime
operations are denied by default.

\paragraph{Filesystem access.}
Filesystem policies constrain authorized objects and access modes, including
read, write, create, unlink, and rename. Enforcement covers both path-based and
file-descriptor-based operations, including file open, path modification,
permission checks on existing descriptors, executable mappings, and controlled
file-descriptor transfer. The trusted runtime resolves and canonicalizes policy
resources before installation. Kernel hooks evaluate bounded predicates over
precomputed mount, inode, file, and namespace identifiers, rather than
performing semantic path resolution in the hot path. Ambiguous cases, such as
paths whose normalization crosses unauthorized mount or symbolic-link
boundaries, are rejected during compilation rather than approximated with
broader access.

\paragraph{Network access.}
Network policies operate on kernel-checkable predicates, including address
family, protocol, destination address or CIDR range, and port. Hostnames and
other semantic endpoint descriptions are resolved before policy installation;
if they cannot be converted into stable predicates, compilation fails closed.
Because network transmissions are irreversible release events, direct external
network access is gated by the ghost's release state. Before semantic approval,
the policy denies direct external transmissions except through trusted mediated
release paths. After approval, connection attempts or datagram transmissions
must still match the installed endpoint predicates.

\paragraph{Process and privilege control.}
Additional hooks restrict process creation, executable loading, capability use,
and executable memory mappings. The runtime clears ambient authority, prevents
privilege-changing transitions, and ensures helper processes remain in the same
invocation cgroup. These checks prevent confined tools from escaping their
declared authority through descendant processes, privileged executables,
unexpected capabilities, or executable code outside the compiled policy.

All hook programs implement default-deny semantics for managed cgroups. If an
operation within the modeled surface is not matched by the active ghost policy,
the kernel rejects the operation.

\subsubsection{Ghost Policies}
\label{sec:design:ebpf:lifecycle}

A \emph{ghost} is the runtime representation of an invocation-specific kernel
policy. Each ghost contains the hot-path enforcement state, the target cgroup
identifier, lease state, release state, and a provenance identifier; the trusted
runtime stores the full provenance record.

Ghosts are created immediately before untrusted execution begins and become
active once their eBPF map entries are installed. During execution, covered
security-sensitive operations issued by the confined process are mediated using
the ghost state associated with the current cgroup.

Ghosts are reclaimed when the invocation completes, is discarded, times out, or
fails. Reclamation removes the corresponding policy entries and releases
execution-specific resources. Each ghost also carries an expiration time checked
by the hook programs; expired ghosts fail closed even if userspace cleanup is
delayed. This lease-based design prevents stale authorization state from
granting privileges after the intended invocation lifetime.

\paragraph{Security properties.}
The enforcement layer maintains three properties:

\begin{enumerate}[leftmargin=1.2em]
\item \textbf{Confinement before untrusted execution.}
The runtime installs the ghost policy and assigns the tool to its cgroup before
transferring control to untrusted tool code.

\item \textbf{Invocation isolation.}
Authorization state is indexed by the invocation's cgroup and is not shared
across independently launched tool invocations.

\item \textbf{Ephemeral privilege.}
Permissions are valid only for the lifetime and release state of the
corresponding ghost and are revoked on completion, discard, timeout, failure,
or lease expiration.
\end{enumerate}

\subsection{Speculative Execution with ShadowFS}
\label{sec:design:speculative}

A fundamental challenge in intent-driven authorization is the latency of
semantic verification. Determining whether a requested permission is consistent
with user intent requires invoking a trusted semantic guard, which typically
incurs hundreds to thousands of milliseconds of latency. A straightforward
design would block tool execution until authorization completes, placing LLM
reasoning directly on the critical path.

\SystemName avoids this bottleneck through \emph{speculative execution under
kernel confinement}. The key observation is that semantic authorization need
not complete before confined tool execution begins; it must complete before the
results of that execution are released. \SystemName therefore allows a tool to
execute under the eBPF/LSM policy derived from SK-IR while treating outputs and
filesystem updates as speculative, and treating externally visible effects as
mediated release requests.

This design decouples \emph{enforcement} from \emph{release synchronization}:
kernel confinement is established before untrusted tool code runs, while
semantic validation proceeds asynchronously in the background.

\subsubsection{Execution Model}
\label{sec:design:speculative:protocol}

Figure~\ref{fig:speculation} illustrates the execution workflow.

After SK-IR compilation succeeds, \SystemName installs the corresponding ghost
policy and launches the tool under kernel confinement. Two activities then
proceed concurrently:

\begin{enumerate}[leftmargin=1.2em]
\item \textbf{Confined speculative execution.}
The tool executes using only the authority granted by the installed eBPF/LSM
policy. Filesystem mutations are redirected to ShadowFS, and tool outputs are
buffered as speculative results rather than immediately returned to the agent.

\item \textbf{Semantic authorization.}
The semantic guard evaluates the task context, accepted permissions, and the
pending release request, including buffered outputs and filesystem or external
effects.
\end{enumerate}

Unlike a blocking authorization design, \SystemName allows confined execution
to overlap with semantic verification. However, authorization synchronization
is required before any speculative effect crosses a release boundary. Release
boundaries include exposing tool output to the agent, committing filesystem
updates to the host environment, issuing externally visible network or API
effects, and allowing a subsequent non-speculative action to depend on
speculative state.

As a result, semantic verification latency can often be absorbed by tool
execution time and by the agent's planning time between release points, reducing
authorization overhead on the critical path without exposing unauthorized
results.

\subsubsection{ShadowFS Overlay}
\label{sec:design:speculative:shadowfs}

ShadowFS is a FUSE-based copy-on-write overlay that isolates speculative
filesystem state from the underlying host filesystem. Each invocation receives
an invocation-local writable overlay. Reads observe committed state and any
explicitly tracked speculative dependencies that the runtime allows the
invocation to consume, while writes, creates, deletes, renames, and metadata
updates are recorded in the overlay rather than applied directly to the host
filesystem.

Because host state is not modified during speculative execution, denial does
not require compensating updates to undo host-visible effects. If semantic
authorization fails, times out, or the invocation crashes, \SystemName discards
the invocation's overlay and buffered outputs and revokes the associated ghost
policy. The unauthorized execution therefore leaves no committed local state.

If authorization succeeds, \SystemName promotes the speculative overlay through
a bounded commit protocol. Updates that can be made visible atomically, such as
single-file replacement via rename, are committed directly. Multi-object
updates are serialized under runtime control and become visible only after the
dependency closure has been authorized. Before promotion, the runtime
revalidates the base objects and policy predicates used during speculation. If
committed host state has changed in a way that invalidates the recorded
dependency or authorization context, promotion fails closed. Operations that
cannot be committed without violating filesystem boundaries, policy
constraints, or crash-safety requirements are rejected rather than approximated
with broader authority.

\subsubsection{Handling Irreversible Operations}
\label{sec:design:speculative:irreversible}

Not all operations can be made safe through filesystem isolation. Externally
visible operations, such as outbound network communication, API invocations,
IPC to host services, or interactions with remote systems, cannot be undone
once released.

\SystemName therefore treats irreversible operations as release boundaries.
Direct external sockets from confined code are denied before release approval.
Tools that require irreversible effects must issue them through a trusted
runtime mediation endpoint, which records the pending effect and releases it
only after approval. If the semantic guard has not yet approved the
corresponding release request, the mediation path withholds the effect and
prevents it from becoming externally observable. Upon approval, the effect is
released subject to both the semantic decision and the installed kernel policy.
Upon denial, timeout, or failure, the pending effect is discarded and the ghost
policy is reclaimed.

This design preserves speculation for reversible local computation while
preventing unauthorized irreversible effects from becoming externally visible.

\subsubsection{Dependency-Aware Consistency}
\label{sec:design:speculative:dependency}

Speculative execution introduces consistency challenges when multiple tool
invocations interact through shared state. A later invocation may read data
produced by an earlier speculative invocation whose semantic authorization has
not yet completed.

To preserve causal consistency, ShadowFS tracks dependencies between
speculative invocations. When an invocation reads data from another uncommitted
overlay or consumes a speculative tool output, the runtime records a dependency
edge. A dependent invocation may continue executing speculatively over this
merged view, but its own outputs and filesystem effects cannot be released
until all upstream dependencies have been authorized.

If an upstream invocation is denied, times out, or fails, \SystemName discards
that invocation's overlay and recursively discards the speculative state of any
dependent invocations. This ensures that committed state and agent-visible
outputs never depend on data produced by an unauthorized execution.

\subsubsection{Security Analysis}
\label{sec:design:speculative:safety}

Speculative execution preserves the security guarantees established by
\SystemName because speculation changes when effects are released, not what
authority the tool receives.

\paragraph{Confinement before untrusted execution.}
The ghost policy is installed and the tool is assigned to its cgroup before
untrusted tool code begins execution. Speculation therefore does not introduce
an unconstrained execution window.

\paragraph{Least-privilege preservation.}
Speculative execution does not expand the authority available to the tool.
Throughout execution, covered security-sensitive operations remain subject to
the kernel policy compiled from SK-IR.

\paragraph{Fail-closed visibility.}
Semantic denial, timeout, crash, or internal failure prevents speculative
effects from crossing release boundaries. Buffered outputs are discarded,
uncommitted overlays are removed, and the corresponding ghost policy is
revoked.

\paragraph{External side-effect safety.}
Irreversible operations are treated as release boundaries and are never
externally issued before semantic authorization succeeds. Unauthorized
invocations therefore cannot produce network-visible, API-visible, or
host-service-visible effects.

The critical observation is that ShadowFS does not rely on undoing
unauthorized host modifications after the fact. Instead, it prevents
speculative effects from becoming observable until semantic authorization
succeeds, and discards them if authorization fails.
